\documentclass[10pt]{article}
\usepackage[margin=0.85in]{geometry}
\usepackage{amsmath,amssymb,amsthm}
\usepackage[hidelinks]{hyperref}
\newtheorem{theorem}{Theorem}
\title{Scalar E-eigenvalues versus sign-pairs: conjecture 00000002844}
\author{earthking11}
\date{October 6, 2026}
\begin{document}
\maketitle
\section*{The precise counting convention}
The source claims that an order-$d$, $n$-dimensional symmetric tensor has
$((d-1)^n-1)/(d-2)$ E-eigenvalues, without identifying opposite scalar
values or imposing a dimension lower bound. We use standard normalized
scalar E-eigenvalues. For a real tensor complexified to $\mathbb C^n$,
these satisfy
\[
 \sum_i x_i^2=1,\quad x\neq0,\qquad
 (Tx^{d-1})_i=\lambda x_i.
\]
The normalization uses the complex-bilinear extension of the Euclidean
quadratic form, not a Hermitian norm. The primary definition calls the
scalar $\lambda$ itself an E-eigenvalue. For odd order,
$(\lambda,x)$ and $(-\lambda,-x)$ are both normalized eigenpairs.
See Sodomaco, Definition 1.1 \cite{sodomaco}.

\begin{theorem}
The one-dimensional order-three symmetric tensor $T_{111}=1$ has exactly
two distinct scalar E-eigenvalues, $1$ and $-1$. The displayed formula
equals $1$, so it does not count these E-eigenvalues as stated.
\end{theorem}
\begin{proof}
There is one tensor coefficient, so all permutations of its three
indices leave it unchanged: it is symmetric. The actual tensor
contraction is $Tx^2=x^2$. The normalized eigenvalue equations reduce to
\[
 x^2=1,\qquad x\neq0,\qquad x^2=\lambda x.
\]
The factorization $(x-1)(x+1)=0$ over $\mathbb C$ implies
$x=1$ or $x=-1$. Substitution then gives $\lambda=1$ or $\lambda=-1$,
respectively. Conversely each of these pairs satisfies every equation.
Thus the actual scalar-value set is exactly $\{1,-1\}$, with cardinality
$2$. At $d=3,n=1$, the asserted formula is
$((3-1)^1-1)/(3-2)=1$, not $2$.
\end{proof}

\section*{What this does and does not refute}
This is a disproof of the source's literal scalar-E-eigenvalue count.
It is not a counterexample to a formula counting values modulo the
identification $\lambda\sim-\lambda$: in this example there is exactly
one sign-pair, agreeing with the displayed number.
The standard odd-order theorem indeed counts sign-pairs, while the
E-characteristic polynomial has twice that pair-count degree
\cite{sodomaco}. Confusing these two quantities is precisely the issue.
The Lean proof establishes the actual scalar-value set and its cardinality;
it does not assert that it formalizes a resultant construction or its
algebraic degree. If the intended conjecture is instead sign-pair counting,
that altered convention must be stated, and this counterexample does not
apply to it. No such identification appears in the source.

\section*{Formal verification}
Lean defines the actual tensor on \texttt{Fin 1}, its symmetry, double-sum
contraction, quadratic unit and nonzero conditions, eigenpair equations,
and scalar-value set with vector witnesses. It proves both directions of
the exact set characterization and then computes its cardinality.
The project builds on Lean 4.33.1 with pinned Mathlib, without unfinished
proofs, native evaluation, extra axioms, or long brute force. Its audited
dependencies are propositional extensionality, classical choice, and
quotient soundness.

\begin{thebibliography}{1}
\bibitem{sodomaco} Luca Sodomaco, \emph{The Product of the Eigenvalues of a
Symmetric Tensor}, arXiv:1802.10173, 2018. Definition 1.1, p. 1;
Theorem 1.3 and the subsequent degree discussion, p. 3.
\url{https://arxiv.org/abs/1802.10173}.
\end{thebibliography}
\end{document}
